\begin{abstract}
3D Gaussian Splatting overfits when only a few views are available, and a common remedy is to supervise rendered depth with a pretrained depth network. Depth Anything~3 (DA3) also returns a per-pixel confidence, which suggests weighting the depth loss so that unreliable predictions count less. We test this idea in a controlled study on top of DropGaussian, with five loss variants that differ only in per-pixel weights, including a control that uses the same weights at shuffled positions. We evaluate on LLFF (3 views, 8 scenes, 3 seeds) and on \emph{Hacettepe Statues}, a new set of \statuesnscenes{} 360$^\circ$ phone captures of outdoor sculptures (12 views) on which nothing was tuned, and we measure the calibration of the confidence and the accuracy of the reconstructed geometry against dense reference reconstructions. A plain Pearson depth loss improves the baseline by \LLFFVoneminusVzerodpsnr{}\,dB PSNR on LLFF and \StatuesVoneminusVzerodpsnr{}\,dB on the statues, helps in \BothVoneminusVzerowins{} scenes and makes geometry more accurate. Confidence weighting adds a small gain on the statues only, and the shuffled control matches it: over all 25 scenes, placing the weights by confidence instead of at random changes PSNR by \BothVtwominusVfourdpsnr{}\,dB ($p=\BothVtwominusVfourp$), and hard masking by confidence loses most of the benefit. This happens although the confidence ranks depth errors better than chance and better than distance alone, and although confidence weighting does give the most accurate geometry of the well-observed region. The confidence is low over whole distant regions, so weighting by it mainly removes supervision where sparse-view reconstruction needs it: coverage of depth supervision matters more than its accuracy.
\end{abstract}
